\clearpage
\VerdierPhysicalPageAnchor{0147}%
\begingroup
\small
\setlength{\parskip}{0pt}%
\setlength{\emergencystretch}{1.5em}%
\setlength{\abovedisplayskip}{3pt plus 1pt minus 1pt}%
\setlength{\belowdisplayskip}{3pt plus 1pt minus 1pt}%

\VerdierSemanticAnchor{ch2-2-3-4-d-complete}
\noindent This proves d); we now prove f).

\VerdierSemanticAnchor{ch2-2-3-4-f}
An object $Z$ of $\mathcal D$ belongs to
$Q^{-1}(\operatorname{Lib}_R(Q))$ if and only if there exist an object $Y$
that is $Q$-free on the right and an isomorphism in
$\mathcal D/\mathcal B$, $Q(Z)\xrightarrow[\sim]{m}Q(Y)$.

\VerdierSemanticAnchor{ch2-2-3-4-f-fraction}
Such an isomorphism has the form $m=Q(s)^{-1}Q(t)$, where
$s\in S_{\mathcal D}(\mathcal B)$ and
$t\in\widehat{S_{\mathcal D}(\mathcal B)}$ ((2.2.2) and (2.2.4)). Hence
in $\mathcal D$ there is a diagram:

\VerdierSemanticAnchor{ch2-2-3-4-f-diagram}
\begin{center}
  \begin{tikzcd}[column sep=4.1em,row sep=2.5em]
    Z \arrow[r,"t"] & X & Y \arrow[l,"s"']
  \end{tikzcd}
  \quad .
\end{center}

\VerdierSemanticAnchor{ch2-2-3-4-f-retraction}
By a), iv), the morphism $s$ admits a retraction
$s'\colon X\longrightarrow Y$ which belongs to
$S_{\mathcal D}(\mathcal B)$ (1.2.6). Hence there exists a morphism
$t'\colon Z\longrightarrow Y$,
$t'\in\widehat{S_{\mathcal D}(\mathcal B)}$.

\VerdierSemanticAnchor{ch2-2-3-4-f-converse}
Conversely, if such a morphism exists, the morphism
$Q(t')\colon Q(Z)\longrightarrow Q(Y)$ is an isomorphism.

\VerdierSemanticAnchor{ch2-2-3-4-g}
It remains to prove g). By a), v), there is a bifunctorial isomorphism:

\VerdierSemanticAnchor{ch2-2-3-4-g-q1-adjunction}
\[
  \operatorname{Hom}_{Q^{-1}(\operatorname{Lib}_R(Q))}(Z,\Psi(W))
  \xrightarrow{\sim}
  \operatorname{Hom}_{\operatorname{Lib}_R(Q)}(Q_1(Z),W),
\]
which shows that $\Psi$ is right adjoint to $Q_1$.

\VerdierSemanticAnchor{ch2-2-3-4-g-inclusion-adjunction}
It follows immediately that, if
$i\colon L_R(Q)\lhook\joinrel\longrightarrow
Q^{-1}(\operatorname{Lib}_R(Q))$ denotes the inclusion functor, there is a
bifunctorial isomorphism:
\[
  \operatorname{Hom}_{Q^{-1}(\operatorname{Lib}_R(Q))}(Z,i(Y))
  \xrightarrow{\sim}
  \operatorname{Hom}_{L_R(Q)}(\Psi Q_1(Z),Y)\, .
\]

\VerdierSemanticAnchor{ch2-2-3-4-g-proof-complete}
This shows that the functor $\mathcal R=\Psi Q_1$ is left adjoint to the
inclusion functor and completes the proof of Proposition (2.3.3).

\begingroup\itshape
\VerdierSemanticAnchor{ch2-2-3-5}
\medskip
\noindent\textbf{Proposition 2.3.5.}\quad
\VerdierSemanticAnchor{ch2-2-3-5-setup}
Let $\mathcal D$ be a triangulated category, let $\mathcal A$ and
$\mathcal B$ be two full triangulated subcategories ($\mathcal B$ strictly
full), and let $S_{\mathcal D}(\mathcal B)$ and
$S_{\mathcal A}(\mathcal A\cap\mathcal B)$ be the multiplicative systems
of $\mathcal D$ and $\mathcal A$ associated respectively with the
categories $\mathcal B$ and $\mathcal A\cap\mathcal B$.

\VerdierSemanticAnchor{ch2-2-3-5-a}
\par\noindent a) The following conditions are equivalent:

\VerdierSemanticAnchor{ch2-2-3-5-a-i}
\par\noindent\hspace{1.2em}i) Every morphism
$X\xrightarrow{s}X'$, with $X\in\operatorname{Ob}(\mathcal A)$ and
$s\in S_{\mathcal D}(\mathcal B)$, fits into a commutative diagram:

\VerdierSemanticAnchor{ch2-2-3-5-a-i-diagram}
\begin{center}
  \begin{tikzcd}[column sep=4.1em,row sep=2.6em]
    X \arrow[r,"s"] \arrow[dr,"{s'}"'] & X' \arrow[d] &
    {s'\in S_{\mathcal A}(\mathcal A\cap\mathcal B)} \\
    & X'' &
  \end{tikzcd}
  \quad .
\end{center}
\endgroup
\endgroup
